\chapter{Dominance et front de Pareto}\label{ch:pareto}
\iterbanner{V1}

\section{Dominance}
Avec deux objectifs à minimiser, deux profils ne sont pas toujours comparables : l'un peut traîner moins et être plus visible. On compare donc composante par composante.
\begin{definition}[ (dominance de Pareto)]
Le profil $a$ \textbf{domine} le profil $b$ si $f_k(a)\le f_k(b)$ pour tous les objectifs $k$, avec au moins une inégalité stricte. Un profil faisable qu'aucun autre profil faisable ne domine est \textbf{non dominé}, ou Pareto-optimal. Leur ensemble est le \textbf{front de Pareto}.
\end{definition}
Le front est la réponse complète d'un problème bi-objectif : il contient tous les compromis rationnels, et le choix entre eux appartient au concepteur, pas à l'algorithme.

\begin{figure}[htbp]
\centering
\begin{tikzpicture}[scale=1.1]
\draw[arr,cmuted] (0,0) -- (7,0) node[below,font=\small,cink] {$f_1$ (traînée)};
\draw[arr,cmuted] (0,0) -- (0,5) node[left,font=\small,cink] {$f_2$ (SER)};
\fill[cblue!8] (3,1.6) rectangle (6.8,4.8);
\node[cblue,font=\small] at (5.2,4.3) {région dominée par $a$};
\foreach \x/\y in {1/4.2,1.6/3,2.3/2.3,3/1.6,4.2/1.05,5.6/0.7}{\fill[cink] (\x,\y) circle (2pt);}
\draw[cink,thick] (1,4.2) -- (1.6,3) -- (2.3,2.3) -- (3,1.6) -- (4.2,1.05) -- (5.6,0.7);
\node[font=\small,anchor=south west] at (3,1.6) {$a$};
\foreach \x/\y in {3.8,3.1/2.7,5/2.4,4.5/3.7,6/1.9,2.6/3.9}{}
\foreach \x/\y in {3.8/3.1,5/2.4,4.5/3.7,6/1.9,2.6/3.9,1.9/4.5}{\fill[cmuted] (\x,\y) circle (1.6pt);}
\node[font=\small,align=center] at (1.6,0.7) {front de\\ Pareto};
\draw[-{Stealth[length=2mm]},cink] (1.6,1.05) -- (2.2,2.1);
\fill[corange] (0.85,0.6) circle (2pt) node[below right,font=\small,corange] {idéal};
\end{tikzpicture}
\caption{Dominance. Les points noirs forment le front : aucun n'est battu sur les deux objectifs à la fois. Le point idéal (meilleure valeur de chaque objectif) n'est en général pas atteignable.}\label{fig:dominance}
\end{figure}

\section{Ramener le problème à un seul objectif : scalarisation}
Beaucoup de solveurs (et tout QUBO) ne savent minimiser qu'\emph{une} fonction. Deux façons classiques de s'y ramener :
\begin{itemize}
\item la \textbf{somme pondérée} $\min\ w\,\tilde f_1+(1-w)\,\tilde f_2$, avec des objectifs normalisés $\tilde f_k\in[0,1]$, et un balayage de $w\in[0,1]$. Chaque $w$ donne un point du front. Mais seuls les points sur l'\emph{enveloppe convexe} du front, dits \textbf{supportés}, peuvent être obtenus ainsi : un point dans un creux du front n'est le minimum d'aucune somme pondérée ;
\item l'\textbf{$\varepsilon$-contrainte} $\min f_1$ sous $f_2\le\varepsilon$, avec un balayage de $\varepsilon$. Elle atteint tous les points du front, convexes ou non, mais ajoute une contrainte d'inégalité, donc des qubits d'écart dans un QUBO.
\end{itemize}

\begin{figure}[htbp]
\centering
\begin{tikzpicture}[scale=1.1]
\draw[arr,cmuted] (0,0) -- (6,0) node[below,font=\small,cink] {$\tilde f_1$};
\draw[arr,cmuted] (0,0) -- (0,4.5) node[left,font=\small,cink] {$\tilde f_2$};
\foreach \x/\y in {0.5/4,1.2/2.4,2.4/2.0,3.3/1.6,4.8/0.5}{\fill[cink] (\x,\y) circle (2pt);}
\draw[cink,thin] (0.5,4) -- (1.2,2.4) -- (2.4,2.0) -- (3.3,1.6) -- (4.8,0.5);
\draw[cblue,thick,dashed] (0.5,4) -- (1.2,2.4) -- (4.8,0.5);
\node[corange,font=\small,anchor=south west] at (2.4,2.0) {non supporté};
\fill[corange] (2.4,2.0) circle (2.4pt); \fill[corange] (3.3,1.6) circle (2.4pt);
\draw[cgreen!60!black,thick] (0.2,2.9) -- (3.2,0.85);
\node[cgreen!60!black,font=\small,anchor=north] at (3.6,0.8) {droite $w\tilde f_1+(1-w)\tilde f_2=\text{cte}$};
\node[cblue,font=\small] at (4.2,2.6) {enveloppe convexe};
\end{tikzpicture}
\caption{La somme pondérée ne trouve que les points de l'enveloppe convexe du front. Les points orange, dans un creux, ne sont le minimum d'aucune combinaison $w\tilde f_1+(1-w)\tilde f_2$.}\label{fig:convexe}
\end{figure}

\section{Mesurer la qualité d'un front approché}
Pour comparer des méthodes, il faut chiffrer la distance entre un front trouvé et le front exact.
\begin{itemize}
\item La \textbf{fraction du front exact trouvée} : proportion des profils Pareto-optimaux qui ont été simulés. C'est le critère le plus strict, et celui de la note d'intention (« atteindre le front exact »).
\item L'\textbf{hypervolume} (HV) : l'aire de la région dominée par le front trouvé et bornée par un point de référence $\vect r$ (figure~\ref{fig:hv}). Plus il est grand, mieux c'est ; il est maximal pour le front exact. On rapporte l'hypervolume trouvé à celui du front exact, sur des objectifs normalisés, avec $\vect r=(1{,}1;\,1{,}1)$.
\item L'\textbf{IGD} (distance générationnelle inversée) : la distance moyenne de chaque point du front exact au point trouvé le plus proche.
\end{itemize}

\begin{figure}[htbp]
\centering
\begin{tikzpicture}[scale=1.1]
\draw[arr,cmuted] (0,0) -- (6,0) node[below,font=\small,cink] {$\tilde f_1$};
\draw[arr,cmuted] (0,0) -- (0,4.5) node[left,font=\small,cink] {$\tilde f_2$};
\fill[cblue!15] (0.6,4) -- (0.6,3.8) -- (1.4,3.8) -- (1.4,2.2) -- (2.6,2.2) -- (2.6,1.4) -- (4.2,1.4) -- (4.2,0.6) -- (5.4,0.6) -- (5.4,4) -- cycle;
\foreach \x/\y in {0.6/3.8,1.4/2.2,2.6/1.4,4.2/0.6}{\fill[cblue] (\x,\y) circle (2pt);}
\fill[corange] (5.4,4) circle (2.4pt) node[above right,font=\small,corange] {référence $\vect r$};
\node[cblue!70!black,font=\small] at (3.8,2.8) {hypervolume};
\end{tikzpicture}
\caption{L'hypervolume est l'aire dominée par le front et bornée par le point de référence. En 2D, c'est une somme de rectangles, après tri selon $f_1$.}\label{fig:hv}
\end{figure}

\begin{codebox}[ : \code{optim/pareto.py}]
\begin{lstlisting}
def hypervolume_2d(F, ref):
    F = F[np.all(F < ref, axis=1)]          # points qui dominent la reference
    F = F[pareto_front(F)]
    F = F[np.argsort(F[:, 0])]              # tri selon f1, f2 decroissant
    hv, prev = 0.0, ref[1]
    for f1, f2 in F:
        hv += (ref[0] - f1) * (prev - f2)   # un rectangle par point
        prev = f2
    return hv
\end{lstlisting}
\end{codebox}

\exercices
\begin{exercice}
Trois profils ont pour objectifs $A=(60;\,-16{,}3)$, $B=(58;\,-16{,}2)$ et $C=(62;\,-16{,}1)$. Lesquels sont non dominés ? Calculer l'hypervolume de $\{A,B\}$ par rapport à $\vect r=(65;\,-16)$.
\end{exercice}
\begin{exercice}
Montrer qu'un point minimisant $w\tilde f_1+(1-w)\tilde f_2$ pour un $w\in\,]0,1[$ est forcément Pareto-optimal.
\end{exercice}

\begin{pistes}
\begin{itemize}
\item Ajouter la scalarisation de Tchebycheff, $\min\max_k w_k|\tilde f_k-z_k^\star|$ : elle atteint les points non supportés sans ajouter de contrainte, mais elle n'est pas quadratique (le max), donc pas directement un QUBO.
\item Utiliser un troisième objectif (le moment de tangage, ou la SER à une autre fréquence) : le tri non dominé et l'hypervolume se généralisent, la visualisation devient plus difficile.
\end{itemize}
\end{pistes}
